% ============================================================
%  Technical Notes: Native Sparse Attention (NSA)
%  Reference: Yuan et al. (DeepSeek, 2025). arXiv:2502.11089
%  Compile: tectonic nsa_notes.tex
% ============================================================

\documentclass[11pt,letterpaper]{article}

\usepackage[T1]{fontenc}
\usepackage[utf8]{inputenc}
\usepackage{geometry}
\geometry{top=1in, bottom=1in, left=1.25in, right=1.25in}

\usepackage{amsmath, amssymb, amsthm}
\usepackage{booktabs}
\usepackage{tabularx}
\usepackage{enumitem}
\usepackage{xcolor}
\usepackage{hyperref}
\usepackage{microtype}
\emergencystretch=2em
\usepackage{parskip}
\usepackage{titlesec}
\usepackage{mdframed}

\hypersetup{colorlinks=true, linkcolor=blue!60!black,
            citecolor=green!50!black, urlcolor=blue!60!black}

\titleformat{\section}{\large\bfseries}{\thesection.}{0.5em}{}
\titleformat{\subsection}{\normalsize\bfseries}{\thesubsection}{0.5em}{}

\title{\textbf{Technical Notes: Native Sparse Attention (NSA)}}
\author{Reference: Yuan et al.\ (DeepSeek, 2025)\\
        \small\texttt{arXiv:2502.11089}}
\date{}

\begin{document}
\maketitle

% ============================================================
\section{Problem Statement}
% ============================================================

Full self-attention scales quadratically with sequence length, making long-context pretraining
computationally prohibitive. Prior sparse attention methods such as BigBird and Longformer
address this through fixed, pre-defined patterns applied after full-attention pretraining. This
post-hoc approach has two consequences. First, training still requires full attention, so
computational savings apply only at inference. Second, because the model is not trained to be
sparse, quality degrades when sparsity is applied. The theoretical speedups in prior methods do
not translate to real hardware gains because the sparse patterns are not aligned with how modern
GPUs load data.

% ============================================================
\section{Contribution}
% ============================================================

NSA (Native Sparse Attention) introduces natively trainable sparse attention in which sparsity
is part of the model architecture from the start of pretraining. The method is designed around
hardware constraints, specifically the memory-bandwidth bottleneck in long-sequence attention on
GPUs. Speedups are demonstrated across all three phases: decoding, forward propagation, and
backward propagation. The model is pretrained on 270 billion tokens at a context length of 64k
tokens.

% ============================================================
\section{Inference Phases and the Prefilling Stage}
% ============================================================

Inference in autoregressive language models consists of two sequential stages: prefilling and
decoding. Understanding their distinct computational profiles is necessary to evaluate where
sparse attention yields the largest gains.

\subsection{Prefilling}

Prefilling is the initial phase in which the model processes the entire input prompt before
generating any new tokens. The model reads all input tokens in parallel and builds the
key-value (KV) cache that decoding will consume. This stage involves two forms of
intensive pre-processing. The first is computing attention maps over the full input sequence,
which has cost $O(n^2)$ in the sequence length $n$. The second is building any auxiliary
indices required by sparse attention, such as block importance rankings for the selection
branch or compressed block representations for the compression branch. Because all input tokens
are available simultaneously, prefilling is compute-bound rather than memory-bandwidth-bound.
For long prompts, it is therefore the phase where sparse attention reduces the most
floating-point operations.

\subsection{Autoregressive Decoding}

Decoding generates tokens one at a time. At each step, the model attends over the KV cache
accumulated from all prior tokens. Because only one new token is produced per step, the
computation is bandwidth-bound: the bottleneck is loading the KV cache from high-bandwidth
memory (HBM) rather than performing matrix multiplications. NSA reduces this bandwidth cost
through GQA-consistent block selection, which loads each relevant KV block once per query group
rather than once per query head (see Section~\ref{sec:gqa}).

% ============================================================
\section{Architecture}
% ============================================================

NSA decomposes attention into three parallel branches that are combined via learned scalar gates.

\subsection{Block Parameters}

Three hyperparameters govern the granularity of each branch. The compression block length $l$
defines how many consecutive tokens are aggregated into a single compressed representation; the
paper sets $l = 32$. The selection block length $l'$ defines the granularity of blocks that the
selection branch retrieves as full, uncompressed tokens; this is set to $l' = 64$. The
compression stride $d$ controls the spacing between adjacent compression blocks; this is set to
$d = 16$, so compression blocks overlap by $l - d = 16$ tokens. Together, these values mean
that the compression branch sees the sequence at a $32\times$ coarser resolution, while the
selection branch operates on blocks twice as large as the compression blocks, which reduces
index-building overhead during prefilling.

\subsection{Compression Branch}

The compression branch produces a coarse-grained representation of the full context. Contiguous
token blocks of length $l = 32$ with stride $d = 16$ are each compressed into a single token
using a learnable MLP with intra-block positional encoding. The model attends over these
compressed block tokens, which reduces the number of keys per query by approximately $32\times$.
The compressed representations also serve as the source of importance scores for the selection
branch: after computing attention between a query and the compressed block tokens, the resulting
attention weights indicate which regions of the sequence are most relevant to that query.

\subsection{Selection Branch}

The selection branch uses the attention scores from the compression branch as importance
estimates. Specifically, for each query $q_t$, the model computes
$\tilde{p}_{\text{slc}} = \mathrm{Softmax}(q_t^\top \tilde{K}_{\text{cmp}})$, then ranks
selection blocks of length $l' = 64$ by their aggregated score and retrieves the top-$k$ blocks
as full, uncompressed tokens. Selection is dynamic: each query retrieves different blocks
depending on content, so the model learns which parts of the context are relevant rather than
following a fixed pattern.

\subsection{Sliding Window Branch}

The sliding window branch attends to a fixed window of the $w$ most recent tokens. It handles
local continuity independently of the other two branches, which prevents gradient interference
between local and long-range learning signals.

\subsection{Gating}

The three branch outputs are combined as:
\[
  o_t \;=\; \sum_{c \,\in\, \{\text{cmp},\,\text{slc},\,\text{win}\}}
            g_t^c \cdot \mathrm{Attn}\!\left(q_t,\,\tilde{K}_t^c,\,\tilde{V}_t^c\right)
\]
where $g_t^c \in [0,1]$ are gate scores produced by a per-branch MLP with sigmoid activation.
The gates allow the model to weight branches differently depending on the query.

% ============================================================
\section{GQA Integration}
\label{sec:gqa}
% ============================================================

NSA is designed for Grouped Query Attention (GQA), in which multiple query heads share a single
KV head. Block selection in the selection branch is made consistent across all query heads within
a GQA group by aggregating their importance scores before ranking:
\[
  \tilde{p}'_{\text{slc}} \;=\; \sum_{h} \tilde{p}_{\text{slc},(h)}
\]
This ensures the same KV blocks are loaded once per group rather than once per head, which
reduces KV-cache bandwidth during decoding. GQA itself reduces the number of KV heads relative
to query heads so that less memory is loaded per decoding step. NSA's group-consistent selection
extends this bandwidth reduction to the selection branch.

% ============================================================
\section{Hardware Motivation and Arithmetic Intensity}
\label{sec:hw}
% ============================================================

\subsection{Arithmetic Intensity}

Arithmetic intensity is the ratio of floating-point operations (FLOPs) to memory accesses
(bytes loaded/stored). It determines which resource is the bottleneck on a given GPU.

\begin{center}\small
\renewcommand{\arraystretch}{1.3}
\begin{tabular}{@{}lll@{}}
\toprule
\textbf{Regime} & \textbf{Intensity vs.\ hardware threshold} & \textbf{Bottleneck} \\
\midrule
Compute-bound  & Intensity \emph{above} critical threshold & GPU FLOP throughput \\
Memory-bound   & Intensity \emph{below} critical threshold & Memory bandwidth (HBM) \\
\bottomrule
\end{tabular}
\end{center}

Every GPU has a critical arithmetic intensity equal to its peak FLOP rate divided by its memory
bandwidth. A workload above this threshold is compute-bound; below it, memory-bound.

\subsection{Phase-Specific Intensity in Language Models}

\begin{center}\small
\renewcommand{\arraystretch}{1.3}
\begin{tabular}{@{}lll@{}}
\toprule
\textbf{Phase} & \textbf{Intensity} & \textbf{Optimisation goal} \\
\midrule
Training (forward + backward) & High --- batched matmuls & Reduce total FLOPs \\
Prefilling                    & High --- parallel sequence processing & Reduce total FLOPs \\
Autoregressive decoding       & Low --- load entire KV cache per token & Reduce memory access \\
\bottomrule
\end{tabular}
\end{center}

Decoding is memory-bandwidth bound because tokens are generated one at a time, each requiring
the full KV cache to be loaded from HBM to SRAM. No batching over the sequence dimension is
possible, so there is little computation to amortise the memory cost.

\subsection{How NSA Achieves Near-Optimal Arithmetic Intensity}

Long-sequence attention is memory-bandwidth bound on modern GPUs. The time to load key-value
tensors from HBM to on-chip SRAM exceeds the time spent on matrix multiplication. NSA reduces
the volume of KV data loaded per query through three mechanisms.

\begin{enumerate}[nosep]
  \item \textbf{Compression branch.} Replaces $l = 32$ raw tokens with one compressed token,
    reducing KV volume by $32\times$ for the global context pass.
  \item \textbf{Selection branch.} Loads only the top-$k$ KV blocks ranked as relevant by the
    query, skipping the rest of the cache entirely.
  \item \textbf{GQA-consistent block selection.} Loads each selected block once per query group
    rather than once per query head, eliminating redundant KV-cache transfers across heads.
\end{enumerate}

Together these keep arithmetic intensity near the hardware-optimal point across all phases:
during training and prefilling (compute-bound), they reduce total FLOPs; during decoding
(memory-bound), they reduce total HBM reads. The paper characterises this as
``arithmetic intensity-balanced'' design.

% ============================================================
\section{Comparison to Prior Methods}
% ============================================================

\begin{table}[h]
\centering
\begin{tabularx}{\textwidth}{lXX}
\toprule
\textbf{Property} & \textbf{BigBird / Longformer} & \textbf{NSA} \\
\midrule
Sparse patterns     & Fixed, pre-defined             & Learned during pretraining \\
Training mode       & Full attention pretraining required & Native sparse pretraining \\
Speedup phases      & Inference only                 & Pretraining + inference \\
Hardware alignment  & Not optimised                  & Explicitly hardware-aligned \\
\bottomrule
\end{tabularx}
\caption{Comparison of NSA against prior sparse attention methods.}
\end{table}

\subsection{Phase-Restricted Sparsity in Older Methods}

A key failure mode of prior work is applying sparsity in only one phase of the model lifecycle,
leaving the other phases at full-attention cost.

\begin{center}\small
\renewcommand{\arraystretch}{1.3}
\begin{tabular}{@{}lll@{}}
\toprule
\textbf{Method} & \textbf{Sparsity applied} & \textbf{Left unaddressed} \\
\midrule
H2O             & Decoding only          & Prefilling, training \\
MInference      & Prefilling only        & Decoding, training \\
Quest, ClusterKV & Inference only        & Training \\
NSA             & Training + prefilling + decoding & Nothing \\
\bottomrule
\end{tabular}
\end{center}

This restriction matters for real workloads. Book summarisation is dominated by prefilling;
long reasoning chains are dominated by decoding. A method that optimises only one phase fails
to accelerate the other. NSA provides hardware-aligned sparse attention computation throughout
the full model lifecycle, which is its central practical advantage over these methods.

% ============================================================
\section{Results}
\label{sec:results}
% ============================================================

\begin{table}[h]
\centering
\begin{tabular}{lcc}
\toprule
\textbf{Benchmark} & \textbf{Full Attention} & \textbf{NSA} \\
\midrule
LongBench (64k avg)          & baseline & $+0.032$ \\
Needle-in-Haystack (64k)     & ---      & Perfect retrieval \\
AIME reasoning (16k)         & 0.092    & 0.146 \\
\bottomrule
\end{tabular}
\caption{NSA versus full attention on long-context benchmarks.}
\end{table}

NSA matches or exceeds full attention on general benchmarks and outperforms it on long-context
retrieval and reasoning tasks.
Notably, NSA achieves \textbf{11$\times$ faster decoding} than full attention, the largest
reported speedup for any phase. This gain is explained directly by the arithmetic intensity
analysis in Section~\ref{sec:hw}: decoding is memory-bandwidth bound, and NSA's three-mechanism
KV reduction (compression, selection, GQA-consistent loading) cuts HBM reads far more
aggressively than any FLOPs reduction could at that phase.

% ============================================================
\section{Sparsity as a Functional Regularizer}
% ============================================================

Beyond efficiency, sparsity in NSA acts as a regularizer on the learning process.
Two mechanisms explain this.

\paragraph{Hierarchical focus.}
The three-branch architecture forces the model to maintain two levels of awareness simultaneously:
the compression branch provides a global, low-frequency scan of the full context, while the
selection branch provides fine-grained access to specific important blocks.
This prevents the model from being overwhelmed by the volume of long-context data.
Without the coarse branch, the model has no global signal when the context is too long for
any individual token to attend to the whole.
Without the fine branch, the model loses precision on tokens that matter.
The hierarchical structure enforces that both are always present.

\paragraph{Suppression of irrelevant context.}
By attending only to selected blocks, NSA prevents the model from incorporating irrelevant
tokens into its representations.
In standard attention over long sequences, weak but non-zero attention weights to distant,
unrelated tokens act as low-level noise.
Sparse selection eliminates this noise structurally rather than relying on the model to learn
to suppress it.
The result is that representations carry more signal from causally relevant tokens, which
strengthens the model's ability to capture long-range logical dependencies and produces more
robust reasoning over extended contexts.

This regularization effect helps explain a counterintuitive result in the NSA paper: NSA
\emph{outperforms} full attention on several long-context benchmarks (see Section~\ref{sec:results})
rather than merely matching it.
The sparse pattern is not a constraint degrading quality; it is a structural prior that
improves it.

% ============================================================
\section{Relation to SparseRenderFormer}
% ============================================================

NSA's three-branch decomposition maps closely to the GP-Sparse module in SparseRenderFormer:

\begin{table}[h]
\centering
\begin{tabularx}{\textwidth}{lX}
\toprule
\textbf{NSA Branch} & \textbf{GP-Sparse Equivalent} \\
\midrule
Compression   & Coarse branch (BVH-pooled block tokens) \\
Selection     & Selected branch (top-$k$ BVH block selection) \\
Sliding window & Local branch (kNN in 3D with normal-coherent mask) \\
Learned gates & Per-head scalar gates (softmax-normalised $\lambda_1, \lambda_2, \lambda_3$) \\
\bottomrule
\end{tabularx}
\caption{Correspondence between NSA branches and GP-Sparse in SparseRenderFormer.}
\end{table}

The primary distinction is that GP-Sparse replaces the position-agnostic sliding window with a
geometry-aware local branch that uses 3D kNN and a normal-coherent mask (suppressing neighbours
where $\mathbf{n}_i \cdot \mathbf{n}_j < -0.5$). This substitution is appropriate for
triangle-mesh rendering, where spatial proximity and surface orientation are more informative
locality signals than token recency.

\end{document}
